\section{Introduction}

Colombia's near-zero aggregate total factor productivity (TFP) growth is not evidence that every economic activity stagnated. The same aggregate result can reflect either similarly weak outcomes across activities or large positive and negative sectoral contributions that offset one another. The distinction matters because the diagnosis changes: a common slowdown calls for examining economy-wide constraints, whereas offsetting contributions require identifying what happened within the activities on each side of the balance.

This paper asks whether Colombia's aggregate TFP stagnation between 2005 and 2024 was a common sectoral outcome or the net result of divergent activity paths. We use the production-account results published by the Departamento Administrativo Nacional de Estad\'istica (DANE) for the total economy and nine broad economic activities. The DANE framework follows the KLEMS growth-accounting approach: output growth is decomposed into the contributions of labor, capital, energy, materials, services, and a TFP residual \citep{DANE2022TFP,OECD2001Productivity,LAKLEMS2021}. The production approach provides a comparable activity history because intermediate inputs enter explicitly.

The paper makes three contributions. First, it converts DANE's annual TFP changes into chained indices for each activity and the total economy. This permits a direct comparison between the levels of 2005 and 2024. Second, it recovers the annual aggregation weights implicit in the published tables and uses them to obtain exact activity contributions to aggregate TFP. The recovery reproduces all 13 published growth-accounting components for the total economy to numerical precision. Third, it uses the same accounting identity to construct a transparent counterfactual. Annual TFP growth is set to zero in the four activities with the weakest long-run performance, while all other input contributions and annual weights remain fixed.

The main result is that aggregate stagnation was real, but it was not a common sectoral pattern. The chained index for the total economy fell by 1.5 percent. Agricultural TFP increased by 27.6 percent, whereas mining TFP fell by 35.1 percent. Social services made the largest positive contribution to aggregate TFP, at 0.148 percentage points per year. Finance and real estate made the largest negative contribution, at $-0.137$ points, followed by mining at $-0.133$ points. Manufacturing and construction each subtracted about 0.081 points per year, despite very different sectoral TFP paths and average weights.

The accounting counterfactual illustrates the aggregate importance of the negative contributions. Production-account output grew by 3.15 percent per year. If annual TFP growth in mining, construction, finance and real estate, and manufacturing had been zero, the same accounting identity yields growth of 3.59 percent per year. The difference is 0.43 percentage points. This is not an estimate of a feasible policy intervention. It does not model changes in factor demand, prices, input use, or resource allocation, and it does not refer to GDP growth.

The rest of the paper is organized as follows. Section~2 describes the DANE data, the KLEMS identity, the long-run chaining procedure, the activity decomposition, and the counterfactual. Section~3 presents the aggregate and sectoral results. Section~4 discusses interpretation, limitations, and policy questions. Section~5 concludes. The appendix documents the weight recovery, additional growth-accounting results, subperiod contributions, and the activity classification.
